\documentclass[11pt,a4paper]{article}
\usepackage[utf8]{inputenc}
\usepackage[vietnamese]{babel}
\usepackage{amsmath,amssymb}
\usepackage{geometry}
\geometry{a4paper, total={180mm,260mm}, left=15mm, right=15mm, top=18mm, bottom=18mm}
\usepackage{tikz}
\usetikzlibrary{shapes.geometric, arrows.meta, positioning, calc}
\usepackage{tcolorbox}
\tcbuselibrary{skins,breakable}
\usepackage{xcolor}
\usepackage{multicol}
\usepackage{array}
\usepackage{fontawesome}
\usepackage{fancyhdr}

\pagestyle{fancy}
\fancyhf{}
\renewcommand{\headrulewidth}{0.5pt}
\renewcommand{\footrulewidth}{0.5pt}
\fancyhead[L]{\small\bfseries\color{teal}\faBook\ SỔ TAY KIẾN THỨC TOÁN 3}
\fancyhead[R]{\small\bfseries\color{teal}Lớp Toán Thầy Lĩnh -- loptoanthaylinh.vn}
\fancyfoot[C]{\small\color{gray}-- Sổ tay kiến thức: Bài 54. Phép cộng trong phạm vi 10 000 --}

\definecolor{primary}{HTML}{005F73}
\definecolor{secondary}{HTML}{0A9396}
\definecolor{accent}{HTML}{EE9B00}
\definecolor{danger}{HTML}{AE2012}
\definecolor{boxbg}{HTML}{F4F6F8}

\begin{document}

% --- BANNER TIÊU ĐỀ ---
\begin{tcolorbox}[
  enhanced,
  colback=white,
  colframe=primary,
  arc=3mm,
  boxrule=1.2pt,
  left=4mm, right=4mm, top=2.5mm, bottom=2.5mm,
  drop shadow=black!10!white
]
  \centering
  {\Large\bfseries\color{primary}\faBookmark\ BÀI 54: PHÉP CỘNG TRONG PHẠM VI 10 000}\\[1.5mm]
  {\small\itshape\color{black!75}Kĩ thuật đặt tính thẳng cột, cộng từ phải sang trái, phép cộng không nhớ và có nhớ không liên tiếp trên 1 trang A4}
\end{tcolorbox}

\vspace{-2mm}


\noindent
\begin{minipage}[t]{0.49\textwidth}\vspace{0pt}
\begin{tcolorbox}[
  enhanced,
  colback=blue!4!white,
  colframe=blue!75!black,
  arc=2mm, boxrule=0.8pt,
  left=2mm, right=2mm, top=1.8mm, bottom=1.8mm,
  title={\footnotesize\bfseries\color{white}\faKey\ 1. Quy trình Đặt tính rồi tính}
]
{\footnotesize
\begin{itemize}\setlength{\itemsep}{1.5pt}\setlength{\leftskip}{-2.5mm}
  \item \textbf{Đặt tính thẳng cột:}
  Viết số hạng này dưới số hạng kia sao cho các chữ số ở \textbf{cùng một hàng thẳng cột} với nhau (đơn vị thẳng đơn vị, chục thẳng chục, trăm thẳng trăm, nghìn thẳng nghìn).
  \item \textbf{Thứ tự cộng:}
  Cộng lần lượt từng hàng từ \textbf{phải sang trái}:
  $$\text{Đơn vị} \to \text{Chục} \to \text{Trăm} \to \text{Nghìn}$$
  \item \textbf{Quy tắc có nhớ:} Nếu tổng ở một hàng $\ge 10$, viết chữ số đơn vị, \textbf{nhớ $1$} sang hàng tiếp theo bên trái.
\end{itemize}
}
\end{tcolorbox}
\end{minipage}
\hfill
\begin{minipage}[t]{0.49\textwidth}\vspace{0pt}
\begin{tcolorbox}[
  enhanced,
  colback=teal!4!white,
  colframe=teal!75!black,
  arc=2mm, boxrule=0.8pt,
  left=2mm, right=2mm, top=1.8mm, bottom=1.8mm,
  title={\footnotesize\bfseries\color{white}\faCalculator\ 2. Tính nhẩm các số tròn nghìn, tròn trăm}
]
{\footnotesize
\begin{itemize}\setlength{\itemsep}{1.5pt}\setlength{\leftskip}{-2.5mm}
  \item \textbf{Cộng số tròn nghìn:}
  $$3000 + 4000 = 7000$$
  (Nhẩm: $3\text{ nghìn} + 4\text{ nghìn} = 7\text{ nghìn}$).
  \item \textbf{Cộng số tròn trăm:}
  $$5200 + 600 = 5800$$
  \item \textbf{Tìm số hạng chưa biết trong phép cộng:}
  $$\text{Số hạng} = \text{Tổng} - \text{Số hạng đã biết}$$
\end{itemize}
}
\end{tcolorbox}
\end{minipage}

\vspace{-2.5mm}
% --- KHỐI 3: SƠ ĐỒ TƯ DUY TIKZ ---
\noindent
\begin{tcolorbox}[
  enhanced,
  colback=orange!4!white,
  colframe=orange!80!black,
  arc=2mm, boxrule=0.8pt,
  left=2mm, right=2mm, top=1.5mm, bottom=1.5mm,
  title={\footnotesize\bfseries\color{white}\faSitemap\ SƠ ĐỒ TƯ DUY: TIẾN TRÌNH CỘNG 4 CHỮ SỐ TỪ PHẢI QUA TRÁI}
]
\centering
\begin{tikzpicture}[
  box/.style={draw, rounded corners=2pt, align=center, font=\scriptsize\bfseries, inner sep=2.5pt, minimum height=0.75cm},
  arr/.style={-{Stealth[length=2mm]}, thick}
]
  \node[box, fill=purple!15, draw=purple!80!black, text width=2.6cm] (c1) at (0, 0) {HÀNG ĐƠN VỊ\\Cộng trước tiên\\$\ge 10 \to$ Nhớ $1$};
  \node[box, fill=teal!15, draw=teal!80!black, text width=2.6cm] (c2) at (3.5, 0) {HÀNG CHỤC\\Cộng thêm nhớ\\$\ge 10 \to$ Nhớ $1$};
  \node[box, fill=blue!12, draw=blue!80!black, text width=2.6cm] (c3) at (7.0, 0) {HÀNG TRĂM\\Cộng thêm nhớ\\$\ge 10 \to$ Nhớ $1$};
  \node[box, fill=orange!15, draw=orange!80!black, text width=2.6cm] (c4) at (10.5, 0) {HÀNG NGHÌN\\Cộng thêm nhớ\\Viết toàn bộ kết quả};

  \draw[arr, purple!80!black] (c1) -- (c2);
  \draw[arr, teal!80!black] (c2) -- (c3);
  \draw[arr, blue!80!black] (c3) -- (c4);
\end{tikzpicture}
\end{tcolorbox}

\vspace{-2.5mm}
% --- KHỐI 4: VÍ DỤ & MẸO TRÁNH BẪY ---
\noindent
\begin{minipage}[t]{0.49\textwidth}\vspace{0pt}
\begin{tcolorbox}[
  enhanced,
  colback=green!4!white,
  colframe=green!60!black,
  arc=2mm, boxrule=0.8pt,
  left=2mm, right=2mm, top=1.8mm, bottom=1.8mm,
  title={\footnotesize\bfseries\color{white}\faCheckCircle\ Ví dụ thực hành chuẩn mực}
]
{\footnotesize
\textbf{1. Phép cộng có nhớ:}\\
$\begin{array}{r@{}l}
  4638 & \\
  +\quad 2517 & \\
  \hline
  \mathbf{7155} &
\end{array}$
\qquad
\begin{minipage}[c]{0.5\linewidth}
-- $8 + 7 = 15$, viết $5$ nhớ $1$.\\
-- $3 + 1 = 4$, thêm $1$ là $5$, viết $5$.\\
-- $6 + 5 = 11$, viết $1$ nhớ $1$.\\
-- $4 + 2 = 6$, thêm $1$ là $7$, viết $7$.
\end{minipage}

\textbf{2. Bài toán thực tế:}\\
Đội Một trồng được $1350$ cây, đội Hai trồng được nhiều hơn đội Một $280$ cây. Đội Hai trồng được:
$$1350 + 280 = \mathbf{1630}\text{ (cây)}$$
}
\end{tcolorbox}
\end{minipage}
\hfill
\begin{minipage}[t]{0.49\textwidth}\vspace{0pt}
\begin{tcolorbox}[
  enhanced,
  colback=red!4!white,
  colframe=red!70!black,
  arc=2mm, boxrule=0.8pt,
  left=2mm, right=2mm, top=1.8mm, bottom=1.8mm,
  title={\footnotesize\bfseries\color{white}\faExclamationTriangle\ Mẹo tránh bẫy trắc nghiệm}
]
{\footnotesize
\begin{itemize}\setlength{\itemsep}{1.5pt}\setlength{\leftskip}{-2.5mm}
  \item \textbf{Bẫy quên cộng số nhớ:}\\
  Khi cộng có nhớ, hãy chấm một dấu nhỏ hoặc ghi số $1$ nhỏ lên trên đầu hàng liền kề bên trái để không bị quên cộng vào!
  \item \textbf{Bẫy đặt tính lệch cột:}\\
  Cộng số có 4 chữ số với số có 3 chữ số (ví dụ $4520 + 315$), phải viết chữ số $3$ thẳng cột với chữ số $5$ (hàng trăm), không viết lệch sang hàng nghìn!
\end{itemize}
}
\end{tcolorbox}
\end{minipage}


\end{document}
